\cvsection{实习经历}

\vspace{-6mm}

\begin{cventries}

    \cventry
      {工程实习生 $\rightarrow$ 研究实习生 | \href{https://dex-robot.com}{DexRobot Inc.}} % Title
      {基于人类示教的高自由度灵巧操作研究} % Project name
      {上海} % Location
      {2024年9月~--~2026年5月} % Date
      {
        \textbf{背景}: 灵巧操作技能学习面临双重瓶颈——遥操作采集灵巧手数据成本高、成功率低、难以规模化; 现有重定向方法多停留在运动学对齐, “形似神不似”, 末端接触难以复现, 人类演示无法直接用于策略训练。围绕硕士课题《基于人类示教的高自由度灵巧操作采集与重定向方法研究》, 独立贯通了从人类演示采集、数据增强到真机部署的完整技术链路。
        \begin{cvitems}
          \item{\textbf{人类演示采集 | IntuitCap 60自由度上肢动捕系统} (一作论文, RCAE 2025已接收): 独立完成机械设计、嵌入式系统与运动学映射的全栈开发——设计3D打印轻量化外骨骼, 集成磁编码器关节角度感知与20自由度液态金属触觉反馈数据手套, 覆盖腰/双臂/腕/双手共60自由度; 基于Socket与CANFD总线实现120Hz采样与7.5ms端到端延迟; 实现运动学建模与映射算法, 以高精度光学动捕为真值验证双臂末端平均误差约10mm; 基于Unity3D数字孪生实现双臂JAKA与DexHand 021实时遥操作}
          \item{\textbf{大规模数据与RL物理增强} (DexCanvas, 共同一作): 贡献基于PPO残差策略的物理回放方法开发, 在Isaac Gym中回放真实示教并提取接触点、接触力与对象六维力/力矩标注, 将70小时真实示教扩展为约7000小时物理增强数据 (覆盖21类操作、30个对象); 统一策略回放成功率87.5\%, 未见YCB物体泛化成功率72.4\%; 接触力标注使下游扩散策略成功率由67.5\%提升至82.2\%}
          \item{\textbf{闭环重定向算法 HOVER} (共同一作, IROS 2025 Workshop Poster): 提出“虚拟操作员+关节残差”两级修正结构, 以近似运动学重定向器提供初始参考, 训练期对人手关键点参考做受限修正, 机器人执行分支输出18维关节残差, 兼顾示教先验保持与执行期误差补偿; DexYCB平均成功率由27.7\%提升至52.0\%, ARCTIC数据集上ADD-AUC较DexMachina最高提升0.131}
          \item{\textbf{策略部署与真机验证}: 完成12主动自由度DexHand 021肌腱驱动灵巧手 (具触觉与力估计) 与JAKA Mini双臂的硬件与控制集成, 实现动态接触条件下的自适应手内操作; 完成网球/玩偶/毛巾抓取搬运、瓶体抓握、双手液体倾倒、测量尺精度跟踪等6类任务, 成功率80\%$\sim$95\%}
        \end{cvitems}
        \vspace{2mm}
      }%
\end{cventries}
